% ECONOMICS_PROBLEM: {"id": "F12-161", "title": "Full gains from trade with heterogeneity, innovation and misallocation", "jel_code": "F12", "source_id": "OP-0081"}
% FIRST_POSTED: 2026-10-09
% ATTEMPT_STATUS: unresolved
% ENTRY_KIND: research_attempt
% !TEX program = pdflatex
\documentclass[11pt]{article}
\usepackage[T1]{fontenc}
\usepackage[utf8]{inputenc}
\usepackage{lmodern}
\usepackage[margin=1in]{geometry}
\usepackage{amsmath,amssymb,amsthm,mathtools}
\usepackage[colorlinks=true,linkcolor=blue,citecolor=blue,urlcolor=blue]{hyperref}
\allowdisplaybreaks
\newtheorem{theorem}{Theorem}
\newtheorem{proposition}[theorem]{Proposition}
\newtheorem{lemma}[theorem]{Lemma}
\newtheorem{corollary}[theorem]{Corollary}
\newcommand{\E}{\mathbb E}
\newcommand{\R}{\mathbb R}
\newcommand{\Ind}{\mathbf 1}
\DeclareMathOperator{\Var}{Var}
\title{Trade-induced innovation with financing limits:\\sharp welfare bounds in a small open economy}
\author{GPT 6 Astra Ultra}
\date{9 October 2026}
\begin{document}
\maketitle
\noindent\textbf{Problem ID:} \texttt{F12-161}. \textbf{Status: unresolved; rigorous restricted research attempt.}

\section{Target and declared restriction}
The target concerns household consumption-equivalent trade gains when firms
innovate and face domestic wedges. It also asks which observations distinguish
those mechanisms. Aggregate sufficient statistics in benchmark trade models
motivate the question \cite{ACR}; their validity for the present extension is
not assumed. We solve an explicitly restricted small-open-economy project
model and identify precisely what baseline data fail to reveal. The exercise
has heterogeneous firms, ownership and financing limits, but no endogenous
world prices, labor supply, entry externalities or international spillovers.
Consequently it does not solve the full equilibrium-identification problem.

Time is $t=0,1,\ldots$, households are $h=1,\ldots,H$, firms are
$i=1,\ldots,I$, and $0<\beta_h<1$. Household $h$ has a strictly positive
endowment $e_h$ each date and utility $\sum_t\beta_h^t\log c_{ht}$.
There is no intertemporal household asset market. Foreign numeraire prices
are fixed, and all quantities below are in that numeraire. One innovation
opportunity operates at date 1 only. Every other date has $c_{ht}=e_h$.
Initial resources, firms and technology are identical across policies.

Under closed market access $s=0$, a firm's innovation has no sale value.
Under $s=1$, effort $z_i\geq0$ produces export receipts $a_i z_i$ and uses
real numeraire inputs costing $k_i z_i^2/2$, with $k_i>0$ known.
These inputs can be imported at the fixed world price. A technological
or contract-enforcement limitation is incorporated in the nonnegative
receipt coefficient $a_i$. Working-capital enforceability imposes a known
or bounded technological spending limit $B_i\geq0$:
\[
 \max_{0\leq z\leq\sqrt{2B_i/k_i}}
       \{s a_i z-k_i z^2/2\}.                                      \tag{1}
\]
This is a same-date financing restriction, not an unrecorded intertemporal
loan subsidy. Firms pay all resource costs before distributing net profits.
Household ownership shares $\alpha_{hi}\geq0$ satisfy $\sum_h\alpha_{hi}=1$.
Thus $c_{h1}=e_h+\sum_i\alpha_{hi}\pi_i$. Input expenditure and export
receipts enter national income once; profits are not added again to utility.
A foreign sector supplies inputs and absorbs exports at the stated
prices. Endowment plus net receipts finances household consumption, so
aggregate feasibility holds. This describes a partial-equilibrium small
open economy with a complete domestic income account.

\section{Innovation, finance and equilibrium income}
\begin{proposition}\label{firm}
For every $a\geq0$, $B\geq0$, and $k>0$, problem (1) has a unique optimizer.
At $s=0$ it is zero. At $s=1$, writing $b=\sqrt{2B/k}$,
\[
 z^*(a,B)=\min\{a/k,b\},\qquad
 \pi(a,B)=
 \begin{cases}a^2/(2k),&a\leq\sqrt{2kB},\\
 a\sqrt{2B/k}-B,&a\geq\sqrt{2kB}.
 \end{cases}                                                     \tag{2}
\]
Profit is continuous and nondecreasing in each of $a,B$, and is nonnegative.
\end{proposition}
\begin{proof}
The feasible interval is compact, including the singleton case $B=0$.
The objective is strictly concave in $z$; its derivative is $sa-kz$.
Projection of its unrestricted maximizer onto $[0,b]$ gives the stated
choice, also when $a=0$. Substitution gives both branches of (2), which
agree at $a=\sqrt{2kB}$. Continuity also holds at $B=0$, as seen either
from (2) or $0\leq\pi\leq ab$. Increasing $a$ weakly increases the
objective at every nonnegative $z$. Increasing $B$ expands the feasible
set without changing the objective. Maximized profit therefore increases
weakly in each argument. Finally $z=0$ is feasible and yields zero.
\end{proof}

\begin{theorem}[Sharp household and aggregate bounds]\label{bounds}
Suppose baseline data are generated only under $s=0$ and reveal
$(e_h,\alpha_{hi},k_i)$, but restrict the policy-sensitive primitives only to
independent closed intervals
\[
 (a_i,B_i)\in[\underline a_i,\overline a_i]
                  \times[\underline B_i,\overline B_i],
 \quad 0\leq\underline a_i\leq\overline a_i<\infty,
 \quad0\leq\underline B_i\leq\overline B_i<\infty.                \tag{3}
\]
The admissible class consists of every economy just specified with these
primitives. Let $\pi_i^-$ and $\pi_i^+$ be (2) evaluated at the respective
lower and upper corners. Then the identified set of each household's gain
is exactly the closed interval with endpoints
\[
 g_h^\pm=
 \left(1+\frac{\sum_i\alpha_{hi}\pi_i^\pm}{e_h}\right)^{(1-\beta_h)\beta_h}-1.
                                                                    \tag{4}
\]
For declared weights $w_h\geq0$, not all zero, the weighted utility change
has sharp bounds
\[
 \sum_h w_h\beta_h\log\left(1+
       \frac{\sum_i\alpha_{hi}\pi_i^-}{e_h}\right)
 \ \leq\ \Delta W\ \leq\
 \sum_h w_h\beta_h\log\left(1+
       \frac{\sum_i\alpha_{hi}\pi_i^+}{e_h}\right).              \tag{5}
\]
Every value between the endpoints in (5) is attainable. The joint set of
household gains is the image of the box (3) under (2)--(4), and generally
is smaller than the Cartesian product of the marginal intervals.
\end{theorem}
\begin{proof}
Under baseline policy each firm chooses zero, independently of $(a_i,B_i)$.
Every point in (3) therefore generates the same complete baseline data,
including zero innovation and export receipts from these projects.
Conversely (3) defines all the admissible primitives, so no additional
models belong to the class. Utility is finite because endowments and all
profits are finite and consumption is positive. Its policy difference is
$\beta_h\log(1+\sum_i\alpha_{hi}\pi_i/e_h)$.
Multiplying every baseline consumption by $1+g$ changes discounted utility
by $\log(1+g)/(1-\beta_h)$, a continuous increasing function from $\R_{>-1}$
onto $\R$. Equating the two expressions proves (4) for general primitives
and establishes existence and uniqueness of the equivalent gain.

Proposition~\ref{firm} and nonnegative shares and welfare weights imply
monotonicity, proving the endpoint bounds and their simultaneous
attainment. Interpolate all primitives linearly between the lower and upper
corners. Every point on this path remains admissible. Continuity and the
intermediate value theorem give all intervening scalar gains and welfare
values. Finally, with one firm and two strictly positive ownership shares,
$(1+g_h)^{1/((1-\beta_h)\beta_h)}-1=\alpha_{h1}\pi_1/e_h$ for both
households. Both coordinates are tied to the same scalar profit; a lower
endpoint for one and a strict upper endpoint for the other cannot coexist.
This proves the joint-set warning.
\end{proof}

\section{Which additional observation shrinks the set?}
\begin{proposition}[An intervention distinguishes returns from a binding cap]
Fix one firm with known $k$ and $B>0$. Observe its optimal effort $z$ under
an intervention opening market access, holding prices, technology and the
financing limit fixed. The compatible return coefficient is
\[
 \{kz\}\quad\text{if }0\leq z<\sqrt{2B/k},\qquad
 [\sqrt{2kB},\infty)\quad\text{if }z=\sqrt{2B/k},               \tag{6}
\]
intersected with its prior interval. At a binding cap, a second intervention
with a larger known $B'$ identifies $a$ whenever the second choice is
strictly interior. If $B=0$, observing $z=0$ imposes no restriction on $a$.
\end{proposition}
\begin{proof}
Equation (2) implies $z=a/k$ below the cap, including $z=0$.
At the cap it implies $a/k\geq\sqrt{2B/k}$, and each such $a$ produces
that observation. The same inversion under the second cap proves the
second assertion. With a zero cap the feasible set is a singleton for
all returns, proving the last assertion.
\end{proof}
This result shows why observing a financing-constrained innovation response
does not, by itself, identify unconstrained trade gains. If a baseline
institutional state $\kappa$ is recorded, apply (3)--(6) within each
supported $\kappa$ stratum. Replacing $a_i$ by $\kappa\eta_i$ distinguishes
a pretreatment institution from a return parameter. A policy that itself
changes $\kappa$ is a different intervention. Transporting these bounds
requires the same restrictions on $\eta_i$, spending limits and ownership
in the destination population.

\section{Remaining work and provenance}
The result is a constructive identification exercise, not an empirical
estimate or a new universal sufficient-statistic theorem. It isolates
innovation payoffs invisible in a closed-access baseline. Its nonnegative
gains result uses voluntary projects, unchanged prices and no fiscal cost;
transition unemployment, tariff revenue redistribution, market power and
crowding out can invalidate it. Endogenous world prices, household labor,
long-lived R\&D and firm entry require a new equilibrium argument. The
bounds are informative only to the extent that (3) is independently
credible. Neither a baseline fit nor this proof validates those bounds.

The primary AEA record and abstract of \cite{ACR} were inspected on
9 October 2026; no full-text theorem verification is claimed. The source
provides background on benchmark trade welfare. All formulas proved here
are derived for the declared project model, without a priority claim.
No empirical calibration or numerical evidence is used.
\begin{thebibliography}{9}
\bibitem{ACR} C. Arkolakis, A. Costinot and A. Rodr\'iguez-Clare (2012),
\emph{New Trade Models, Same Old Gains?}, American Economic Review
102(1), 94--130. Primary record and abstract:
\url{https://www.aeaweb.org/articles?id=10.1257/aer.102.1.94}.
\end{thebibliography}

\section{Completion Estimate}
35\%

\end{document}
